%% ============================================================================
%%  WWW '27 (The ACM Web Conference 2027) submission -- main driver file
%%
%%  Build:   latexmk -pdf main.tex        (or `make` in this directory)
%%  Clean:   latexmk -C                   (or `make clean`)
%%
%%  Review submission  : \documentclass[sigconf, anonymous, review]{acmart}
%%  Camera-ready       : \documentclass[sigconf]{acmart}  (flip \camerareadyfalse)
%%
%%  Page budget (CFP):  8 pages main body, self-contained.
%%                      References + optional appendix may extend to 12 total.
%%                      Short paper: 4 pages INCLUDING references + appendix.
%% ============================================================================

%% ---- Submission mode toggle -------------------------------------------------
\newif\ifcameraready
\camerareadyfalse        % <- set \camerareadytrue for the final version
%% `make camera` passes \def\forcecamera{} on the command line instead.
\ifdefined\forcecamera \camerareadytrue \fi

\ifcameraready
  \documentclass[sigconf]{acmart}
\else
  \documentclass[sigconf, anonymous, review]{acmart}
\fi

%% ---- Packages and project macros -------------------------------------------
%% ============================================================================
%%  preamble.tex -- packages and project-wide macros
%%  acmart already loads: amsmath, amssymb, hyperref, graphicx, xcolor,
%%  booktabs, natbib, microtype, caption, float, geometry.  Do NOT reload them.
%% ============================================================================

%% ==== Tables & floats ========================================================
\usepackage{multirow}
\usepackage{tabularx}
\usepackage{array}
\usepackage{subcaption}        % \begin{subfigure} ... (compatible with acmart)
\usepackage{siunitx}           % \num{}, \SI{}, S-columns for aligned decimals
\sisetup{detect-all, group-separator={,}, group-minimum-digits=4}

%% ==== Algorithms & code ======================================================
\usepackage[ruled,linesnumbered,vlined]{algorithm2e}
\usepackage{listings}
\lstset{basicstyle=\ttfamily\footnotesize, breaklines=true, columns=fullflexible,
        frame=single, framerule=0.3pt, xleftmargin=0.5em, numbers=none}

%% ==== Figures ================================================================
\usepackage{tikz}
\usetikzlibrary{positioning,arrows.meta,calc,fit,shapes.misc}
\usepackage{pgfplots}
\pgfplotsset{compat=1.18}
\usepgfplotslibrary{groupplots}
\graphicspath{{figures/}}

%% ==== Lists ==================================================================
\usepackage[inline]{enumitem}
\setlist{nosep, leftmargin=*}

%% ==== Cross-referencing (load LAST) ==========================================
\usepackage{balance}           % \balance before \bibliography for even columns
\usepackage[capitalise,noabbrev,nameinlink]{cleveref}
\crefname{section}{Section}{Sections}
\crefname{algorithm}{Algorithm}{Algorithms}
\crefname{appendix}{Appendix}{Appendices}

%% ==== Editorial macros (disabled automatically in camera-ready) =============
\ifcameraready
  \newcommand{\todo}[1]{}
  \newcommand{\note}[1]{}
  \newcommand{\cut}[1]{}
\else
  \newcommand{\todo}[1]{\textcolor{red}{\textbf{[TODO: #1]}}}
  \newcommand{\note}[1]{\textcolor{blue}{\textit{[#1]}}}
  \newcommand{\cut}[1]{\textcolor{gray}{#1}}   % text under consideration for removal
\fi

%% ==== Domain macros ==========================================================
%% Edit these once; use them everywhere so renames are painless.
\newcommand{\sys}{\textsc{TaxLex}\xspace}                 % TODO: system / artefact name
\newcommand{\corpus}{\textsc{CryptoDef}\xspace}            % TODO: corpus name
\newcommand{\ie}{i.e.,\xspace}
\newcommand{\eg}{e.g.,\xspace}
\newcommand{\etal}{et al.\xspace}
\newcommand{\cf}{cf.\xspace}
\newcommand{\vs}{vs.\xspace}
\newcommand{\term}[1]{\textit{#1}}                         % a term being defined
\newcommand{\ttype}[1]{\texttt{#1}}                        % a term *type* / class label
\newcommand{\src}[1]{\textsf{#1}}                          % a source instrument (MiCA, FATF, ...)
\usepackage{xspace}

%% ==== Math ===================================================================
\newcommand{\R}{\mathbb{R}}
\newcommand{\E}{\mathbb{E}}
\newcommand{\argmax}{\operatorname*{arg\,max}}
\newcommand{\argmin}{\operatorname*{arg\,min}}
\newcommand{\set}[1]{\{#1\}}
\newcommand{\norm}[1]{\left\lVert #1 \right\rVert}
\DeclareMathOperator{\entropy}{H}
\DeclareMathOperator{\KL}{KL}

%% ==== Space-saving knobs (use sparingly; reviewers notice squeezed papers) --
%\setlength{\textfloatsep}{8pt plus 2pt minus 2pt}
%\setlength{\abovecaptionskip}{4pt}


%% ---- ACM metadata -----------------------------------------------------------
%% Rights/DOI/ISBN are filled in by ACM after acceptance; placeholders below
%% keep the class quiet during review.
\copyrightyear{2027}
\acmYear{2027}
\setcopyright{rightsretained}
\acmConference[WWW '27]{Proceedings of the ACM Web Conference 2027}{May 10--14, 2027}{Dublin, Ireland}
\acmBooktitle{Proceedings of the ACM Web Conference 2027 (WWW '27), May 10--14, 2027, Dublin, Ireland}
\acmDOI{10.1145/XXXXXXX.XXXXXXX}
\acmISBN{978-x-xxxx-xxxx-x/27/05}
\acmSubmissionID{000}     % OpenReview paper number (shown in review mode)

%% Hide ACM reference block + CCS/keywords footer while under review to save
%% space; they are required in the camera-ready.
\ifcameraready\else
  \settopmatter{printacmref=false}
  \renewcommand\footnotetextcopyrightpermission[1]{}
\fi

\begin{document}

%% ---- Title ------------------------------------------------------------------
%% TODO: final title. Keep it specific; name the artefact (taxonomy), the
%% method (LLM-based extraction) and the domain (crypto-assets / DLT).
\title[LLM-Grounded Crypto-Asset Taxonomy]{Grounding a Crypto-Asset Taxonomy in Law and Literature:\\
LLM-Based Term-Type Extraction from Regulatory and Academic Corpora}

%% ---- Authors (hidden by the `anonymous` option) ----------------------------
\author{Anonymous Author(s)}
\affiliation{%
  \institution{Anonymous Institution}
  \city{City}
  \country{Country}}
\email{anonymous@example.org}

%% Short author list for running heads.
\renewcommand{\shortauthors}{Anon. et al.}

%% ---- Abstract ---------------------------------------------------------------
\begin{abstract}
Crypto-assets are issued, traded and described on the Web, and the vocabulary
that governs them is published as Web documents by regulators, standard-setters
and researchers who do not share a dictionary: the same token is a crypto-asset
under MiCA, a digital asset under US law and a virtual asset under FATF, and
each definition discriminates along dimensions the text never names. Manual
taxonomies capture one perspective and go stale, while LLM taxonomy and
ontology-learning methods assume a fixed label space and do not trace their
output to a legal source. In this paper we
propose \sys, a framework that induces the classification dimensions, which we
call term types, from the definitions themselves. \sys (1) builds a
definitions-first corpus from 24 legal instruments and academic taxonomies,
(2) extracts term types in a single pass over the whole corpus held in a cached
context, requiring a document, section and verbatim quote for every dimension
and value, (3) verifies every quote against the source, (4) orders the term
types into an is-a graph by asking, for each one, for a cited probability
distribution over the more general dimensions it specialises, consolidated
greedily into a directed acyclic graph, and (5) uses the graph as the label
space for Web content and grows it only where tagged content raises its
information gain. On a corpus spanning the EU, the UK, the
US and six international bodies, \sys produces 93 term types with 317 values and
480 citations, 97.9\% of which verify, records a cross-source disagreement for 69
of them, and links them by 279 weighted is-a edges under one root per category,
for a few dollars of API calls. Two maintenance iterations over \AllItems\ DLT-Corpus papers and patents
verify 96--97\% of tag evidence and grow the graph by \AllAddedTypes\
dimensions and \AllAddedValues\ values for \$\AllCost, and the dimensions
discovered first are among the most used next. \todo{One sentence of comparative evidence once RQ2--RQ4 are run.} The
corpus, taxonomy and code are available at an anonymised repository.\footnote{\todo{anonymised URL}}

\end{abstract}

%% ---- CCS concepts & keywords ------------------------------------------------
%% Generate the CCS block at https://dl.acm.org/ccs and paste it here.
\begin{CCSXML}
<ccs2012>
   <concept>
       <concept_id>10002951.10003227.10003351</concept_id>
       <concept_desc>Information systems~Data mining</concept_desc>
       <concept_significance>500</concept_significance>
   </concept>
   <concept>
       <concept_id>10010147.10010178.10010179</concept_id>
       <concept_desc>Computing methodologies~Natural language processing</concept_desc>
       <concept_significance>500</concept_significance>
   </concept>
   <concept>
       <concept_id>10002951.10003260.10003309</concept_id>
       <concept_desc>Information systems~Web data description languages</concept_desc>
       <concept_significance>300</concept_significance>
   </concept>
   <concept>
       <concept_id>10003456.10003462</concept_id>
       <concept_desc>Social and professional topics~Computing / technology policy</concept_desc>
       <concept_significance>300</concept_significance>
   </concept>
</ccs2012>
\end{CCSXML}

\ccsdesc[500]{Information systems~Data mining}
\ccsdesc[500]{Computing methodologies~Natural language processing}
\ccsdesc[300]{Information systems~Web data description languages}
\ccsdesc[300]{Social and professional topics~Computing / technology policy}

\keywords{taxonomy induction, large language models, crypto-assets, distributed ledger technology, regulation, knowledge organisation, ontology learning}

\maketitle

%% ============================================================================
%%  MAIN BODY  (must fit in 8 pages, self-contained)
%% ============================================================================
\section{Introduction}
\label{sec:intro}

A taxonomy is a hierarchy of classes that organises a domain's vocabulary so
that machines and people can classify, browse and retrieve its objects
consistently. On the Web, taxonomies drive product catalogues, search
verticals and content moderation, and the community has studied how to expand
and complete them from Web text for years~\cite{shen2020taxoexpan,
wang2021hef,jiang2022taxoenrich,arous2023taxocomplete,shi2024icon}. Crypto-assets
are a domain in which the vocabulary itself is contested. Tokens are issued by
smart contracts, traded on exchanges and decentralised venues, and described in
whitepapers, explorers and social media, and the documents that say what a
token \emph{is} are themselves Web publications, among them the EU's Markets
in Crypto-Assets regulation (MiCA), the US GENIUS and CLARITY Acts, the UK's
cryptoasset regulations, and the guidance of the Financial Action Task Force
(FATF), the Basel Committee, the FSB and IOSCO, alongside academic
taxonomies~\cite{tasca2019taxonomy,zhang2026cryptotaxonomy}.

\paragraph{Relevance to the Web.}
This work falls under the Semantics and Knowledge track. The research challenge
is Web-specific because the definitions that govern crypto-assets are
published, versioned and consumed as Web documents by heterogeneous bodies,
no shared vocabulary links them, and the assets they describe exist only as records on
Web-accessible ledgers. The result is a machine-readable, provenance-tracked
vocabulary that block explorers, exchanges and compliance services can use to
classify tokens consistently across jurisdictions and to regenerate the
classification when an instrument changes.

\begin{figure}[t]
  \centering
  %% Figure 1: one real term, four sources, one induced term type (column width).
\begin{tikzpicture}[
  font=\footnotesize,
  src/.style={draw=gray!70, fill=gray!8, rounded corners=2pt, text width=2.75cm, inner sep=2.5pt, align=left},
  tt/.style={draw=blue!60!black, fill=blue!6, rounded corners=2pt, text width=3.25cm, inner sep=3pt, align=left},
  val/.style={font=\scriptsize, inner sep=1pt},
  arr/.style={-{Latex[length=1.6mm]}, gray!80, thin},
  node distance=2mm and 5mm]
  \node[src] (mica) {\textbf{MiCA} Art.~3(6--7)\\[1pt]\scriptsize ``e-money token \ldots\ referencing the value of \emph{one official currency}''; ``asset-referenced token \ldots\ referencing \emph{another value or right} or a combination''};
  \node[src, below=of mica] (genius) {\textbf{GENIUS Act} \S2(22)\\[1pt]\scriptsize ``maintain \ldots\ a stable value relative to \ldots\ a \emph{fixed amount of monetary value}''};
  \node[src, below=of genius] (uk) {\textbf{UK SI 2026/102} art.~88G\\[1pt]\scriptsize ``a stable value in relation to \emph{a particular fiat currency}''};
  \node[src, below=of uk] (basel) {\textbf{Basel SCO60.6}\\[1pt]\scriptsize Group~1b: ``cryptoassets with effective \emph{stabilisation mechanisms}''; Group~2: ``stablecoins and unbacked cryptoassets''};
  \node[tt, right=of $(genius.east)!0.5!(uk.east)$] (tt) {\textbf{Term type: reference asset}\\[1pt]\scriptsize\emph{What does the asset reference to maintain a stable value?}\\[3pt]
    \scriptsize
    \textbullet~single official currency (EU)\\
    \textbullet~basket or other assets (EU)\\
    \textbullet~fiat currency (US, UK)\\
    \textbullet~commodity (Basel)\\
    \textbullet~other crypto-assets / none (Basel)};
  \draw[arr] (mica.east) -- (tt.west|-mica.east) -- ([yshift=4mm]tt.west);
  \draw[arr] (genius.east) -- ([yshift=1mm]tt.west);
  \draw[arr] (uk.east) -- ([yshift=-1mm]tt.west);
  \draw[arr] (basel.east) -- (tt.west|-basel.east) -- ([yshift=-4mm]tt.west);
  \node[font=\scriptsize\itshape, gray!90, above=1mm of mica.north west, anchor=south west] {term: \emph{stablecoin}};
\end{tikzpicture}

  \caption{One term, four sources, one induced term type. Each instrument
  defines \term{stablecoin} by what the asset references, but the admissible
  references differ by jurisdiction. \sys makes the dimension explicit, keeps
  one value per jurisdictional position, and cites the section each value comes
  from.}
  \label{fig:one}
\end{figure}

These sources disagree in a specific way. \Cref{fig:one} shows the term
\term{stablecoin}. \src{MiCA} splits it into e-money tokens, which reference one
official currency, and asset-referenced tokens, which reference anything else.
The \src{GENIUS Act} and the UK regulations admit only a fiat currency as the
reference. The Basel Committee admits a commodity and treats algorithmic
stabilisation as a different category. None of the texts states that ``what the
asset references'' is a classification dimension, yet each uses it to draw a
legal or prudential boundary. We call such a dimension a \emph{term type}, a
question along which definitions discriminate, together with the set of values
the sources admit. The reference asset is one of many; the issuer's
identifiability, the fungibility of the asset, the custody model and the
consensus mechanism are others, and they recur across instruments under
different names.

Existing approaches do not recover term types from such a corpus. First,
manually built taxonomies of crypto-assets and of blockchain
architectures~\cite{tasca2019taxonomy,zhang2026cryptotaxonomy} encode one
perspective each, a technical or an investor's one, and cannot be regenerated
when an instrument changes. Second, LLM methods for taxonomy
completion~\cite{shi2024icon}, hierarchical
classification~\cite{zhang2025teleclass} and tagging~\cite{wan2024tntllm,
li2023taggpt} start from a label space that already exists, and ontology
learning with LLMs~\cite{babaeigiglou2023llms4ol} is evaluated on general-domain
benchmarks where the gold hierarchy is known. Third, none of these methods ties
an induced class to the passage that licenses it, which in a legal domain is the
difference between a taxonomy a compliance team can audit and one it cannot.

In this paper we propose \sys, a framework that induces term types from the
definitions in a corpus of legal instruments and academic taxonomies and
attaches provenance to every node. \sys builds a definitions-first corpus, in
which each document is reduced to the sections that define or classify; it
extracts term types in one pass over the whole corpus held in a cached context,
so that the model can recognise the same dimension under different names across
sources; it requires a document index, an exact section label and a verbatim
quote for every term type and every value; and it verifies each quote against
the corpus, repairing wrong section labels and rejecting quotes it cannot find.
It then orders the catalogue. For each term type it asks for a cited probability
distribution over the more general term types it specialises, and consolidates
the weighted edges greedily into a directed acyclic graph (DAG), so that a term
type can sit under more than one parent. The output is an ontology of term types
with values, aliases, jurisdictions, recorded conflicts and weighted is-a edges,
every element of which cites its source.

\smallskip\noindent\textbf{Contributions.}
\begin{itemize}
  \item We formalise term type induction, in which a corpus of definitional
    text yields the classification dimensions it uses, with values and
    provenance, without a seed taxonomy or gold labels (\cref{sec:problem}).
  \item We propose \sys, which combines a definitions-first corpus, single-pass
    extraction over a cached whole-corpus context with mandatory citations, a
    citation verifier, and a per-node is-a query consolidated into a DAG with
    metrics for its later maintenance (\cref{sec:method}).
  \item We release \corpus, a definitions-first corpus of 24 documents across
    four jurisdictions, and the resulting ontology of 93 term types with 317
    values, 480 citations and 279 weighted is-a edges
    (\cref{sec:exp:rq1,sec:exp:rq2}).
  \item We evaluate faithfulness and coverage on \corpus: 97.9\% of term-type
    citations verify, every document contributes to at least one term type, 69
    term types record a cross-source conflict, and the hierarchy resolves
    \BBits\ bits per root-to-leaf descent under one root per category. \todo{Add the headline comparative
    number from RQ2--RQ4.} Code and data are at an anonymised
    repository.\footnote{\todo{anonymised URL}}
\end{itemize}

\section{Related Work}
\label{sec:related}

%% Target: 0.75-1 page. Group by theme, end each paragraph with the gap we fill.
%% Use third person for own prior work (double-blind).

\paragraph{Crypto-asset and DLT taxonomies.}
\todo{Component-level blockchain taxonomy~\cite{tasca2019taxonomy};
asset-level taxonomy for investors and regulators~\cite{zhang2026cryptotaxonomy};
ISO/TC~307 vocabulary; regulatory classifications (MiCA token classes, Basel
Group~1/2, FATF VA/VASP). Gap: hand-built, single-perspective, no provenance
across sources.}

\paragraph{LLMs for text classification and tagging.}
\todo{Zero-shot classifiers~\cite{wang2023zeroshot}, clue-and-reasoning
prompting~\cite{sun2023textclassllm}, instruction-tuned
classification~\cite{yu2023llmclassification}, extreme multi-label
in-context learning~\cite{zhu2024icxml}, LLM-based tagging systems
\cite{li2023taggpt,wu2025llm4tag}, label-taxonomy generation at scale
\cite{wan2024tntllm}, and topic discovery via entropy-based clustering
\cite{entropyclustering2024}. Gap: fixed or generated label spaces are not
grounded in authoritative definitions.}

\paragraph{Retrieval-augmented and long-context grounding.}
\todo{RAG~\cite{lewis2020rag}, hierarchical retrieval~\cite{sarthi2024raptor},
RAG with ontologies~\cite{ragontologies2024}, and cost/benefit of orchestrating
multiple LLM calls~\cite{orchestration2025}. Position our cached whole-corpus
grounding against chunk retrieval.}

\paragraph{LLMs for ontology learning and knowledge graphs.}
\todo{LLMs4OL~\cite{babaeigiglou2023llms4ol}; linking crypto attribution tags
to KG entities~\cite{cryptotags2024}; domain corpora for DLT
\cite{hernandez2026dltcorpus}. Gap: no work induces term \emph{types} from legal
text with jurisdiction-aware reconciliation.}

\section{Problem Formulation}
\label{sec:problem}

%% Target: 0.5-0.75 page. Define every symbol once; reviewers reward precision.

\subsection{Terminology}
\todo{Define \term{term}, \term{definition}, \term{term type} (a classification
dimension along which a definition discriminates, \eg \ttype{legal-status},
\ttype{technical-standard}, \ttype{economic-function}, \ttype{issuer-type}),
\term{source instrument} and \term{jurisdiction}.}

\subsection{Corpus Model}
Let $\mathcal{D} = \set{d_1,\dots,d_N}$ be a set of source documents, each
partitioned into labelled sections $d_i = \set{s_{i1},\dots,s_{in_i}}$ with
metadata (type, jurisdiction, year).
\todo{State the two document classes: \emph{legal instruments} and
\emph{taxonomy frameworks}.}

\subsection{Task Definition}
\todo{Formalise: given a term $t$ and $\mathcal{D}$, return a set of term types
$\mathcal{Y}_t = \set{(y, \mathcal{P}_y)}$ where each $y$ carries provenance
$\mathcal{P}_y \subseteq \bigcup_i d_i$ (document $\times$ section citations).
Then define the reconciliation task: merge $\set{\mathcal{Y}_t}_{t\in\mathcal{T}}$
into a taxonomy $\mathcal{G}=(V,E)$ with jurisdiction-indexed nodes.}

\subsection{Desiderata}
\begin{enumerate*}[label=(D\arabic*)]
  \item \emph{faithfulness} -- every term type must be traceable to a cited
    section;
  \item \emph{coverage} -- all instruments that define $t$ are represented;
  \item \emph{consistency} -- equivalent types across sources are merged, conflicting
    ones are preserved with their jurisdiction;
  \item \emph{cost} -- extraction scales to the full corpus within a stated token budget.
\end{enumerate*}

\section{Term Type Generation}
\label{sec:method}

We build the taxonomy in stages. The first stage, which this section describes,
asks a single question of the sources. Along which dimensions do they classify
crypto-assets, the systems that record them, the actors that handle them, the
activities those actors perform, and the obligations attached to each? We call
these dimensions \emph{term types}. Later stages place terms under the term
types and link the term types to one another (\cref{sec:method:hierarchy}).
A term type exists in our output only if the model can quote a passage that
uses it as a criterion, and we check every quote against the corpus after the
run.

\begin{figure*}[t]
  \centering
  \todo{Overview figure: PDFs $\rightarrow$ definitions-first extraction
  $\rightarrow$ XML corpus behind a prompt cache $\rightarrow$ one extraction
  call $\rightarrow$ JSON term types $\rightarrow$ citation verifier
  $\rightarrow$ verified catalogue.}
  %\includegraphics[width=\textwidth]{overview.pdf}
  \caption{Overview of the term type generation stage of \sys. \todo{caption}}
  \label{fig:overview}
\end{figure*}

\subsection{A Definitions-First Corpus}
\label{sec:method:corpus}

We do not give the model whole statutes. A regulation such as \src{MiCA}
defines its vocabulary in a few articles and spends the rest of its length on
obligations that refer back to those definitions. The classification criteria
we are after live in the definitions articles, the glossaries, the
classification conditions and the descriptions of business models. We take
those sections and nothing else. \Cref{tab:corpus} lists the result: 21 legal
instruments from the EU, the UK, the US and six international bodies, plus two
academic taxonomies of blockchains~\cite{tasca2019taxonomy} and of
crypto-assets~\cite{zhang2026cryptotaxonomy} and one paper describing a
domain corpus~\cite{hernandez2026dltcorpus}, with reference lists and large
tables removed.

\begin{table}[t]
  \centering
  \caption{\corpus: the 24 source documents, the sections we kept from each,
  their length, and the number of Task A term types that cite them.}
  \label{tab:corpus}
  \scriptsize
  %% Generated by paper/scripts/task_a_tables.py -- do not edit by hand.
\setlength{\tabcolsep}{3.5pt}
\begin{tabular}{@{}rllrrr@{}}
    \toprule
    \# & Jur. & Source & Sect. & Words\,(k) & Types \\
    \midrule
    1 & Acad. & Tasca \& Tessone taxonomy & 2 & 16.5 & 21 \\
    2 & Acad. & Zhang et al.\ taxonomy & 2 & 9.1 & 24 \\
    3 & Acad. & DLT-Corpus paper & 1 & 6.2 & 1 \\
    4 & EU & MiCA, Reg.\ 2023/1114 & 1 & 2.7 & 25 \\
    5 & EU & AMLR, Reg.\ 2024/1624 & 1 & 4.4 & 11 \\
    6 & EU & Transfer of Funds, Reg.\ 2023/1113 & 1 & 1.0 & 11 \\
    7 & EU & DORA, Reg.\ 2022/2554 & 1 & 2.3 & 5 \\
    8 & UK & FSMA (Cryptoassets) Regs 2026 & 2 & 4.6 & 24 \\
    9 & UK & CARF Regs 2025 & 1 & 0.6 & 3 \\
    10 & US & CLARITY Act (Senate-reported) & 2 & 1.4 & 13 \\
    11 & US & CLARITY Act (House-passed) & 1 & 4.8 & 20 \\
    12 & US & FIT21, H.R.\ 4763 & 1 & 4.7 & 20 \\
    13 & US & GENIUS Act, Pub.\ L.\ 119-27 & 1 & 1.7 & 20 \\
    14 & US & Executive Order 14178 & 1 & 0.1 & 3 \\
    15 & US & PWG Digital Asset Markets Report & 1 & 12.2 & 22 \\
    16 & US & FinCEN CVC Guidance 2019 & 3 & 7.2 & 17 \\
    17 & Intl & FATF VA/VASP Guidance 2021 & 2 & 9.1 & 27 \\
    18 & Intl & FATF Targeted Update 2025 & 1 & 1.7 & 10 \\
    19 & Intl & FATF Targeted Update 2026 & 1 & 2.6 & 12 \\
    20 & Intl & Basel SCO60 & 2 & 3.1 & 15 \\
    21 & Intl & FSB Recommendations 2023 & 1 & 0.5 & 14 \\
    22 & Intl & IOSCO Recommendations 2023 & 1 & 0.7 & 11 \\
    23 & Intl & UN Cybercrime Convention & 1 & 0.7 & 7 \\
    24 & Intl & Budapest Convention & 1 & 0.2 & 3 \\
    \midrule
    \multicolumn{3}{@{}l}{Total} & 32 & 97.8 & 93 \\
    \bottomrule
\end{tabular}

\end{table}

Each PDF is converted with \texttt{pdftotext} and cut by line ranges or by
regular-expression markers that we specify once per document. The US bills
need the layout-preserving mode because their two-column line numbering
scrambles the plain output, and we drop page footers such as the bill number
that would otherwise leak into every section. The Senate-reported print of the
\src{CLARITY Act} reproduces the House-passed text before the Senate
substitute, so we take Title~I once, from the House file, and only the
substitute's definitions and its section on network tokens from the Senate
file. The corpus is one XML-tagged text file. Every document carries an integer
index and metadata (type, jurisdiction, title, year) and is split into
\texttt{<section label="...">} blocks, so a citation can name a document and a
section unambiguously. The file holds 32 sections and 97.8 thousand words,
which the model's tokeniser counts as 235,298 tokens.

\subsection{The Extraction Task}
\label{sec:method:task}

The system prompt defines the objects we want. A \emph{term} is a word or
phrase the corpus defines or classifies, such as \term{crypto-asset} or
\term{permissioned ledger}. A \emph{term type} is a question along which the
corpus's definitions discriminate between things, such as \ttype{reference
asset} or \ttype{consensus mechanism}, and it has \emph{values}, the categories
a thing can take on that dimension. A term type is \emph{grounded} when at
least one passage uses it as a criterion or defines several terms that differ
along it. The model is told that the corpus is its only source of truth and
that outside knowledge may be used to recognise what the corpus says, never to
add to it.

Three grounding rules follow. Every term type and every value must carry at
least one citation of the form \{document index, section label, quote\}; the
quote must be verbatim and at most thirty words; a candidate that cannot be
grounded must be omitted. Four quality rules shape the output: one question per
term type, with differently named versions of the same question merged and the
other names kept as aliases; jurisdictional disagreement about values or
thresholds recorded in a notes field with citations for each side; preference
for dimensions that recur across sources or that a source uses to draw a legal,
prudential or technical boundary, while still keeping an unusual dimension that
a single source grounds clearly; and JSON only.

The user prompt states the scope, which covers assets, DLT systems, actors,
activities, and obligations or risks when the sources themselves use them as
criteria, and gives the JSON schema. Each term type has an identifier, a name,
the category it applies to, the question it answers, a short description, its
aliases, its values with their own definitions and citations, its own
citations, the documents, jurisdictions and document classes that ground it,
and the notes field. The model is asked to order term types by category and
then by the number of documents that ground them, and to aim for completeness.

Our first run cited section labels from memory and got them wrong in
9\% of citations, with 65 term types. We therefore append to the user prompt a
\emph{manifest}: one line per document listing its index, title and every exact
section label, with the instruction that no other labels exist. The manifest is
generated from the corpus at run time, so it cannot drift from the file the
model sees.

\subsection{One Call over a Cached Corpus}
\label{sec:method:run}

We place the whole corpus in the system prompt, after the stable instructions
and behind a prompt-cache breakpoint with a one-hour lifetime; the task prompt
is the user message. Every later task, and every prompt iteration, then pays
only for the small uncached part. This choice is also a methodological one.
Quality rule four asks the model to recognise that \src{MiCA}'s ``referencing
the value of one official currency'' and the \src{GENIUS Act}'s ``fixed amount
of monetary value'' answer the same question. A retriever would have to
anticipate that two passages that share few words define the same dimension;
a model that sees all 24 documents at once does not.

A single call to \texttt{claude-haiku-5-5} at high effort produced the output
we analyse. We disable adaptive thinking because on a first attempt the model
spent more than 100,000 thinking tokens on this corpus without emitting any
text, whereas without thinking it answers directly. We stream the response, allow up
to 120,000 output tokens, and parse the JSON; if a response is cut off, the
parser salvages every complete term type and flags the result as truncated.
The run reported here ended normally after 302 seconds with 85,729 output
tokens, 2,542 uncached input tokens for the task prompt and manifest, and
235,298 input tokens read from the cache.

\subsection{Citation Verification}
\label{sec:method:verify}

A quote that the model labels verbatim may still be invented, compressed, or
attached to the wrong section, so we check all of them. For each citation we
look up the document and section, normalise both the quote and the section text
to lower-case alphanumerics, and accept the quote if it is a substring. If it is
not, we split the quote into 12-character shingles and accept it when at least
80\% of them occur inside one window of the section twice the quote's length.
Localising the matches tolerates footnote markers, hyphenation and the spacing
artefacts of PDF extraction, and rejects a paraphrase stitched from phrases
that occur in different places. When a quote fails under the cited label but passes under another label of
the same document, we \emph{relabel} it, correcting the section, keeping the
model's label alongside, and counting the citation as verified. Everything else is rejected and listed in a report.
Every citation in the published output carries a \texttt{verified} flag, and
the rejection rate is our measure of desideratum~(D1), faithfulness.
Verification establishes that the quote exists where the model says it does;
it does not establish that the quote supports the dimension. We return to this
limit in \cref{sec:exp:rq1}.

\subsection{From Term Types to a Hierarchy}
\label{sec:method:hierarchy}
\todo{Task B: for each term type, a probability distribution over more general
term types in the same category (plus \emph{none}), with rationale and
citations, asked against the same cached corpus plus the Task A catalogue as a
second cached block. Describe once the Task B run is reviewed.}

\section{Experiments}
\label{sec:exp}

We answer six research questions.
\textbf{RQ1}: How many classification dimensions does the corpus support, how
widely is each shared across sources, and how faithful are the citations?
\textbf{RQ2}: What hierarchy do the term types form, how confident and how
grounded are its edges, and how does it change as weak edges are hidden?
\textbf{RQ3}: Does the DAG serve as a label space for Web content, and does
the content reveal where it must grow?
\textbf{RQ4}: How does the induced ontology compare with expert-built
taxonomies and with LLM baselines that are given the same corpus?
\textbf{RQ5}: How much do the manifest, the grounding strategy and the model
matter?
\textbf{RQ6}: Does the ontology improve a downstream classification of
Web-sourced token descriptions?
RQ1 to RQ3 and the manifest part of RQ5 are answered below; the remaining parts
are framed with the tables they will fill.

\subsection{Experimental Setup}
\label{sec:exp:setup}

\textbf{Datasets.} \corpus is described in \cref{tab:corpus}. No standard
benchmark exists for term type induction, so \corpus is both the input and, via
its citations, the ground against which faithfulness is checked. For RQ3 the
content is DLT-Corpus~\cite{hernandez2026dltcorpus}: from its first shard of
scientific papers (37,440 in the collection) and of USPTO patents (49,023) we
keep English items whose title or abstract names a DLT term, which 40\% of
papers and 49\% of patents do, and sample disjoint batches of 150 papers and
150 patents; an item is the title and abstract capped at 450 words. \todo{Held-out set and expert gold set for RQ4 and RQ6; contamination note.}

\textbf{Baselines.} \todo{Expert taxonomies; zero-shot LLM without the corpus; chunked RAG and
RAPTOR over it; TnT-LLM and LLMs4OL-style term typing; same corpus, model and
schema for all.}

\textbf{Metrics.} For faithfulness we report the share of citations whose quote
the verifier finds in the cited document (exact or relabelled), the share
relabelled, and the rejection rate; smaller rejection is better. For coverage we
report the number of term types, values and aliases, the number of documents
grounding each term type, and the number of term types with a recorded
conflict. For the hierarchy we report edges, roots, leaves, depth, bits per descent,
expected depth, AUC and effective leaves as defined in
\cref{sec:method:hierarchy}. \todo{For RQ4 add precision and recall of term
types against the expert gold set, and coverage of each expert taxonomy's
dimensions; for RQ6 add macro-F$_1$ of the downstream classifier.}

\textbf{Implementation.} Both tasks use \texttt{claude-haiku-5-5} called
through the Anthropic Messages API at effort \texttt{high} with adaptive
thinking disabled; the corpus sits behind a one-hour prompt-cache breakpoint.
Task~A is one streamed call with an output limit of 120,000 tokens; Task~B is
93 calls with $k{=}6$, four in parallel, each limited to 4,000 output tokens.
Task~C runs at effort \texttt{medium} with six items in flight; the discovery
and consolidation calls use the defaults of \cref{sec:method:maintain}. The verifier uses 12-character
shingles and an 80\% threshold. All runs are API calls from a laptop; no GPU is
involved. We report a single run. \todo{Add a second run with the same prompt
and report agreement of term types and values, and a run with
\texttt{claude-sonnet-5-5}, before submission.}

\subsection{RQ1: Term Types Generated and Their Faithfulness}
\label{sec:exp:rq1}

\begin{table}[t]
  \centering
  \caption{Term types by category. Docs is the mean number of source documents
  grounding a term type; Single counts term types grounded in one document;
  Confl.\ counts term types whose notes record a disagreement between sources;
  Verif.\ is the share of citations whose quote the verifier found in the cited
  document, exact or relabelled. Single run.}
  \label{tab:taska}
  \scriptsize
  %% Generated by paper/scripts/task_a_tables.py -- do not edit by hand.
\setlength{\tabcolsep}{3.5pt}
\begin{tabular}{@{}lrrrrrrr@{}}
    \toprule
    Category & Types & Values & Docs & Single & Confl. & Cites & Verif.\,\% \\
    \midrule
    Asset & 21 & 68 & 4.5 & 7 & 21 & 101 & 100.0 \\
    Technology & 24 & 82 & 3.2 & 5 & 16 & 123 & 97.6 \\
    Actor & 19 & 67 & 4.1 & 0 & 12 & 103 & 95.1 \\
    Activity & 10 & 43 & 3.9 & 1 & 8 & 63 & 100.0 \\
    Obligation or risk & 19 & 57 & 2.7 & 7 & 12 & 90 & 97.8 \\
    \midrule
    All & 93 & 317 & 3.6 & 20 & 69 & 480 & 97.9 \\
    \bottomrule
\end{tabular}

\end{table}


\textbf{What the corpus supports.} The run produced 93 term types with 317
values and 212 aliases (\cref{tab:taska}). We make the
following observations. (1) Technology dimensions are the largest group,
because the blockchain taxonomy of Tasca and Tessone contributes architectural
dimensions that no legal instrument uses, such as the fee regime and the
interoperability of a ledger. (2) Asset and actor dimensions are the most
widely shared. An asset term type is grounded in 4.5 documents on average, and
the status of an actor as a service provider is grounded in 13 of the 24
documents. (3) Twenty term types rest on a single document, mostly distinctions the
academic taxonomies draw and the law does not yet make, such as the mechanism
by which yield reaches a holder. Every document is cited by at least one term
type; the FATF guidance of 2021 is cited most, by 27, and the DLT-Corpus paper
least, by one; its sections describe data sources and metrics rather than
classification criteria.


Sixty-nine term types carry a note recording a disagreement between sources.
\Cref{tab:case} (\cref{app:results}) shows the reference-asset dimension of
\cref{fig:one} as returned: five values, seven grounding documents across four
jurisdictions, and a note recording the jurisdictional split. The
fungibility dimension likewise records that fungibility is an element of the UK
and US definitions of the asset class but not of the \src{FATF} or \src{MiCA}
definitions. These notes are the raw material for the jurisdiction-indexed
nodes of \cref{sec:method:hierarchy}.


\textbf{Faithfulness.} Of the 480 citations, 463 matched the cited section
verbatim under normalisation, 7 were relabelled, and 10 were rejected, a
rejection rate of 2.1\% (\cref{tab:taska}, last column).1\% for actor dimensions to 100\% for asset and
activity dimensions. One cited label did not exist, a subsection heading inside
the FinCEN guidance that the model used in place of the enclosing section's
label, and six named a real but wrong section of the right document. The verifier located the quote in each of these
seven cases. \Cref{tab:errors} in \cref{app:results} lists the ten rejections with a
diagnosis. Two
quotes share no five-word sequence with the cited section. Three compress a
definition, two drop the enumerated items of a statutory list, as in
``'blockchain' means any technology (i) where data is shared across a
network'', and three reproduce most of a sentence and then rephrase its end.
These quotes point at the right passage but are not verbatim, so we keep them in
the output as unverified rather than repairing them. Every term type keeps at
least one verified citation, but five values were supported only by a rejected
quote and are therefore unverified.

Verification is a necessary check, not a sufficient one. A verbatim quote can
still fail to support the dimension it is attached to, and our verifier does not
read for relevance. Whether the 470 verified citations do support their term
types is the question the expert annotation planned in \cref{sec:exp:setup} is
designed to answer, and we have not run it yet.

\subsection{RQ2: The Is-a Hierarchy}
\label{sec:exp:rq2}

\begin{table}[t]
  \centering
  \caption{The is-a DAG by category after consolidation, no weight threshold.
  Bits/desc.\ is the information gain of a random root-to-leaf descent;
  E[depth] its expected length; AUC the area under the reach-a-leaf curve (1 when
  every root is a leaf); Eff.\,lv.\ is the effective number of leaves
  $2^{\entropy(r)}$ for the leaf-reach distribution $r$. Single run.}
  \label{tab:taskb}
  \scriptsize
  %% Generated by paper/scripts/task_a_tables.py -- do not edit by hand.
\setlength{\tabcolsep}{2pt}
\begin{tabular}{@{}lrrrrrrrrr@{}}
    \toprule
    Category & Nodes & Edges & Roots & Leaves & Depth & Bits/desc. & E[depth] & AUC & Eff.\,lv. \\
    \midrule
    Asset & 21 & 68 & 1 & 5 & 9 & 4.17 & 2.54 & 0.773 & 3.2 \\
    Technology & 24 & 75 & 1 & 4 & 14 & 5.47 & 3.94 & 0.754 & 2.6 \\
    Actor & 19 & 58 & 1 & 3 & 12 & 6.02 & 3.30 & 0.767 & 2.6 \\
    Activity & 10 & 31 & 1 & 2 & 7 & 3.88 & 2.32 & 0.741 & 2.0 \\
    Obligation or risk & 19 & 47 & 1 & 4 & 12 & 4.05 & 1.93 & 0.881 & 4.0 \\
    \midrule
    All & 93 & 279 & 5 & 18 & 14 & 4.72 & 2.81 & 0.835 & 14.0 \\
    \bottomrule
\end{tabular}

\end{table}


\textbf{Structure.} The 93 calls returned \BWeighted\ weighted candidate edges,
of which consolidation kept \BEdges\ and dropped \BDropped\ for closing a cycle,
with \BTies\ exact tie (\cref{tab:taskb}). We make the following observations.
(1) Every category closes under a single root, the crypto-asset definitional
criteria for assets, the governance model for technologies, service-provider
status for actors, the service type for activities, and AML applicability for
obligations; the model was free to return \emph{none} for any node and did so
with median mass 0.13. (2) The DAG is deep and multi-parented: a node keeps
\BMeanParents\ parents on average and up to \BMaxParents, the technology category
reaches depth 14, and a random descent resolves \BBits\ bits before ending at
one of \BEffLeaves\ effective leaves. (3) In the asset category (\cref{fig:tree}, \cref{app:results}) the reference-asset dimension
specialises the legal classification and carries the stablecoin class,
redemption mechanism and tokenisation form beneath it, but second-choice
parents carry nearly as much mass: the model gave
\emph{legal classification} and \emph{stablecoin regulatory class} the same
weight of 0.32 as parents of \emph{reference asset}, and \BMutual\ pairs of
nodes were each proposed as the other's parent. The weights are coarse, with a
median top weight of \BMedianTop, so we read them as rankings rather than
calibrated probabilities.


\begin{figure*}[t]
  \centering
  %% Figure: hierarchy metrics when edges below a weight are hidden (text width, 1x3).
%% Legends sit above each panel so they never cover the curves.
\begin{tikzpicture}
  \begin{groupplot}[
    group style={group size=3 by 1, horizontal sep=15mm},
    width=0.33\textwidth, height=3.6cm,
    xmin=0, xmax=1, xtick={0,0.2,0.4,0.6,0.8,1},
    tick label style={font=\scriptsize}, label style={font=\scriptsize},
    ylabel near ticks, xlabel near ticks,
    axis line style={gray!60}, ymajorgrids, grid style={gray!25},
    legend style={font=\scriptsize, draw=none, fill=none, at={(0.5,1.02)}, anchor=south, legend columns=-1,
                  /tikz/every even column/.append style={column sep=5pt}},
    every axis plot/.append style={mark size=1.4pt, line width=0.9pt}]
    \nextgroupplot[ylabel={count}, xlabel={(a) minimum edge weight}, ymin=0]
      \addplot[blue!60!black, mark=*, mark repeat=10] table[x=thr, y=edges] {tables/spectrum.dat};
      \addplot[orange!80!black, mark=square*, mark repeat=10] table[x=thr, y=roots] {tables/spectrum.dat};
      \addplot[green!50!black, mark=triangle*, mark repeat=10] table[x=thr, y=leaves] {tables/spectrum.dat};
      \legend{kept edges, roots, leaves}
    \nextgroupplot[ylabel={bits per descent}, xlabel={(b) minimum edge weight}, ymin=0]
      \addplot[blue!60!black, mark=*, mark repeat=10] table[x=thr, y=bits] {tables/spectrum.dat};
    \nextgroupplot[ylabel={fraction}, xlabel={(c) minimum edge weight}, ymin=0, ymax=1.05, ytick={0,0.25,0.5,0.75,1}]
      \addplot[blue!60!black, mark=*, mark repeat=10] table[x=thr, y=auc] {tables/spectrum.dat};
      \addplot[orange!80!black, mark=square*, mark repeat=10] table[x=thr, y expr=\thisrow{effleaves}/93] {tables/spectrum.dat};
      \legend{AUC, effective leaves / 93}
  \end{groupplot}
\end{tikzpicture}

  \caption{Sensitivity to the display threshold: the sub-DAG keeps only edges
  with weight at or above the x value. (a) kept edges, roots and leaves; (b)
  information gain per descent; (c) AUC and effective leaves as a fraction of
  the 93 nodes. Hiding edges below 0.4 keeps 64 edges and leaves 30 roots.}
  \label{fig:spectrum}
\end{figure*}

\textbf{Sensitivity to the threshold.} \Cref{fig:spectrum} sweeps a display
threshold over the consolidated DAG. Hiding edges below 0.2 halves the edge
count and roughly halves the bits per descent; at 0.4 a third of the nodes are
roots and a descent resolves half a bit; above 0.6 the graph is a forest of
near-singletons. The full DAG carries the most discrimination; a thresholded
one is easier to read, and the curves price each cut.

\textbf{Faithfulness of edges.} The entries carry \BCites\ citations, of which
\BVerified\ (\BVerifiedPct\%) verify, \BRelabelled\ after relabelling, and
\BRejected\ are rejected, a markedly higher rejection rate than in Task~A. Only
\BEntriesCited\ of the \BWeighted\ entries cite anything, which the prompt
allows for low-probability entries, and \BTopNoCite\ top entries cite nothing,
which it does not. Rejected quotes are mostly short phrases cited from
long sections, such as ``Wrapped tokens''. Verification again checks
existence, not that the quote evidences the relation.
\todo{Expert judgement of a sample of 100 edges on a 1--3 scale with agreement,
following the protocol planned for \cref{sec:exp:rq4}.}

\subsection{RQ3: The DAG as a Label Space}
\label{sec:exp:rq3}

\begin{table*}[t]
  \centering
  \caption{Two maintenance iterations. Each tags a disjoint batch with the
  catalogue version in force, measures the DAG on it, discovers candidates at
  the target nodes and consolidates the accepted ones into the next version.
  Evid.\ is the share of evidence spans found verbatim in the item; Nodes used
  counts nodes that received at least one item; USD and minutes cover tagging,
  discovery and consolidation at list price with eight workers.}
  \label{tab:evolution}
  \scriptsize
  %% Generated by paper/scripts/task_a_tables.py -- do not edit by hand.
\setlength{\tabcolsep}{3pt}
\begin{tabular}{@{}llrrrrrrlllr@{}}
    \toprule
    It. & Cat. & Items & Tags/item & Evid.\,\% & Nodes used & Collapsed & Targets & Cand.\ (acc.) & Produced: nodes/edges/values & USD & Min. \\
    \midrule
    1 & v0 & 300 & 3.55 & 96.3 & 80/93 & 3 & 15 & 32 (8) & v1: 99/285/351 & 0.15 & 2 \\
    2 & v1 & 8,646 & 4.06 & 96.8 & 98/99 & 11 & 32 & 81 (21) & v2: 107/293/396 & 4.29 & 34 \\
    \bottomrule
\end{tabular}

\end{table*}

\begin{table}[t]
  \centering
  \caption{Objectives of each catalogue version from its is-a weights (prior)
  and from the tags of the batch classified with it.}
  \label{tab:objectives}
  \scriptsize
  %% Generated by paper/scripts/task_a_tables.py -- do not edit by hand.
\setlength{\tabcolsep}{3pt}
\begin{tabular}{@{}lrrrrrrrrrrrr@{}}
    \toprule
    Objective & v0 prior & iter.\ 1 tags & v1 prior & iter.\ 2 tags & v2 prior & iter.\ 3 tags & v3 prior & iter.\ 4 tags & v4 prior & iter.\ 5 tags & v5 prior & iter.\ 6 tags \\
    \midrule
    Bits per descent & 4.72 & 6.92 & 4.82 & 6.96 & 4.84 & 7.02 & 4.88 & 7.22 & 5.01 & 7.27 & 5.02 & 7.06 \\
    AUC of reach-a-leaf & 0.835 & 0.794 & 0.843 & 0.771 & 0.851 & 0.778 & 0.843 & 0.780 & 0.833 & 0.773 & 0.833 & 0.773 \\
    Expected depth & 2.81 & 3.37 & 2.70 & 3.70 & 2.74 & 3.83 & 2.86 & 3.80 & 3.00 & 3.90 & 3.00 & 3.90 \\
    Effective leaves & 14.0 & 19.7 & 18.5 & 21.3 & 22.5 & 21.7 & 23.1 & 24.7 & 25.7 & 24.8 & 26.1 & 22.9 \\
    Effective leaves / nodes & 0.150 & 0.212 & 0.187 & 0.215 & 0.210 & 0.203 & 0.204 & 0.218 & 0.207 & 0.200 & 0.204 & 0.179 \\
    \bottomrule
\end{tabular}

\end{table}

\textbf{Tagging.} \Cref{tab:evolution} summarises both iterations. The pilot
batch of \CItems\ items received \CTags\ tags, \CTagsPerItem\ per item, and the
second batch of \TItems\ items, a tenth of the on-topic collection, received
\TTags\ tags, \TTagsPerItem\ per item; \CEvidencePct\% and \TEvidencePct\% of
the evidence spans were found verbatim in the item, and \CRejected\ and
\TRejected\ tags were rejected for unknown identifiers. We make the following
observations. (1) The content exercises the technical half of the DAG: in the
second batch \TCovTech\ items carry a technology tag and \TCovAsset\ an asset
tag, against \TCovActor\ for actors, \TCovActivity\ for activities and
\TCovOblig\ for obligations. The legal nodes are dormant on the pilot batch and collapsed on a single
value on the larger one; testing those branches needs regulatory content. (2) The dimensions discovered in
iteration 1 are among the most used in iteration 2: ledger purpose was tagged
on \TLedgerPurpose\ items with a value entropy of \TEntLedgerPurpose\ bits,
smart-contract usage domain on \TUsageDomain, identity binding on \TIdentity\
and the privacy-mechanism family on \TPrivacy\ with \TEntPrivacy\ bits, so
they discriminate the content rather than absorbing it. (3) Tagging cost
\$\TCostPerItem\ per item at \TMeanLat\ seconds per request; the whole
on-topic collection of about \OnTopicK\ thousand documents would cost about
\$\ProjCost\ and \ProjHours\ hours at the same concurrency.

\textbf{Measurement.} \Cref{tab:objectives} compares, for each version, the
objectives from the model's is-a weights with those measured from the tags. A
content-driven descent resolves about seven bits against under five from the
prior and ends at more effective leaves: the content goes deeper and spreads
wider than the prior weights suggested. Collapsed nodes grow with the batch,
\TCollapsed\ in iteration 2, led by the DLT data structure with 96\% of 2,377
items on ``blockchain''; they are also the largest residual pools, which is
where discovery is pointed.

\begin{table}[t]
  \centering
  \caption{Term types added by the two iterations, with the node they were
  discovered under, the number of shown residual items supporting them and the
  increase in branching gain at that node in bits. Values added to existing
  nodes (\EAddedValues\ and \TAddedValues) are not listed.}
  \label{tab:accepted}
  \scriptsize
  %% Generated by paper/scripts/task_a_tables.py -- do not edit by hand.
\setlength{\tabcolsep}{3pt}
\begin{tabularx}{\columnwidth}{@{}lX>{\raggedright\arraybackslash}p{0.3\columnwidth}rr@{}}
    \toprule
    It. & New term type & Under & Items & Gain \\
    \midrule
    1 & Privacy mechanism family & Privacy and anonymity features & 8 & 0.83 \\
    1 & Smart contract deployment context & Smart contract / script language expressiveness & 6 & 0.66 \\
    1 & Smart contract usage domain & Smart contract / distributed ledger application scope & 8 & 0.46 \\
    1 & Identity binding mechanism & Identity layer and onboarding & 5 & 0.34 \\
    1 & Ledger purpose & Recording and cryptographic linkage & 6 & 0.21 \\
    1 & Ledger partitioning or layering topology & DLT data structure (blockchain vs non-chain DLT) & 5 & 0.18 \\
    2 & Cross-chain transfer mechanism & Wrapped / cross-chain representation & 8 & 0.74 \\
    2 & Key custody and control model & Wallet storage mode (hot or cold) & 7 & 0.67 \\
    2 & Finality mechanism & Finality of agreement & 5 & 0.55 \\
    2 & Key isolation from network connectivity & Hosted versus unhosted (self-hosted) wallet & 5 & 0.55 \\
    2 & Payment settlement mechanism & Payment service and funds transfer & 5 & 0.38 \\
    2 & Key splitting scheme & Permitted control over client assets (unilateral vs multi-party control) & 8 & 0.32 \\
    2 & Expressiveness claim basis & Smart contract / script language expressiveness & 5 & 0.25 \\
    2 & Key management control model & Crypto-asset service type & 6 & 0.11 \\
    3 & Zero-knowledge proof system variant & Privacy mechanism family & 6 & 0.61 \\
    3 & Key share distribution topology & Key management control model & 5 & 0.59 \\
    3 & Fiat-rail integration & Payment settlement mechanism & 5 & 0.55 \\
    3 & Key recovery mechanism & Hosted versus unhosted (self-hosted) wallet & 6 & 0.50 \\
    3 & Script verification approach & Smart contract / script language expressiveness & 5 & 0.26 \\
    3 & Backing of token & Tokenisation form & 5 & 0.14 \\
    4 & Deanonymisation / tracing approach & Anonymity-enhancing crypto-asset status & 6 & 0.75 \\
    4 & Backing mechanism & Backing of token & 5 & 0.72 \\
    4 & Settlement finality timing & Payment settlement mechanism & 5 & 0.42 \\
    4 & Settlement counterparty & Fiat-rail integration & 5 & 0.65 \\
    4 & Physical transfer mechanism across air gap & Key isolation from network connectivity & 5 & 0.57 \\
    4 & Validator selection mechanism & Validator operation capacity (who may run a node) & 5 & 0.55 \\
    4 & Sharding partition basis & Ledger partitioning or layering topology & 5 & 0.54 \\
    4 & Wallet interface or delivery form & Hosted versus unhosted (self-hosted) wallet & 6 & 0.44 \\
    4 & Share recovery mechanism & Key splitting scheme & 5 & 0.44 \\
    4 & Centralisation locus & Centralisation of control & 6 & 0.38 \\
    4 & Cross-chain interoperability mechanism & Wrapped / cross-chain representation & 7 & 0.28 \\
    5 & Key isolation mechanism & Key isolation from network connectivity & 7 & 0.59 \\
    5 & Key distribution method across isolation boundary & Key isolation from network connectivity & 5 & 0.48 \\
    5 & Finality guarantee mode & Finality of agreement & 5 & 0.41 \\
    5 & Key share cryptographic primitive & Key splitting scheme & 5 & 0.30 \\
    6 & ZKP application role & Zero-knowledge proof system variant & 5 & 0.86 \\
    6 & Key storage medium or environment & Key custody and control model & 5 & 0.79 \\
    6 & Parent-child coupling mechanism & Ledger partitioning or layering topology & 6 & 0.58 \\
    6 & Finality guarantee source & Finality of agreement & 5 & 0.47 \\
    \bottomrule
\end{tabularx}

\end{table}

\textbf{Discovery and evolution.} Iteration 1 found \MTargets\ targets and
accepted \MAccepted\ of \MCandidates\ candidates; iteration 2 found \TTargets\
targets and accepted \TAccepted\ of \TCandidates, producing version 2 with
\VTwoNodes\ term types, \VTwoEdges\ edges and \VTwoValues\ values
(\cref{tab:accepted}). The additions are dimensions a legal corpus would not
contain, such as the cryptographic family behind a privacy guarantee or the
custody model of a key, and several iteration-2 values refine iteration-1
nodes. Two failures are visible. Two application-domain dimensions proposed
under two smart-contract nodes in iteration 1 ask the same question and were
not merged by the consolidation call. And with raw support counts, capped by
the forty items shown, only 10 of 76 candidates cleared the gain test against
residual pools of thousands in iteration 2; extrapolating support to the
residual is the rule we report, and the rerun accepted 21 of 81.

The tags have no gold labels yet, so verified evidence is the only quality
signal we report; an evidence span can exist without the tag being right.
\todo{Expert review of a slice of the tags for precision and recall, and a
regulatory content batch for the legal branches.}

\subsection{Planned Evaluations (RQ4--RQ6)}
\label{sec:exp:rq4}
\label{sec:exp:rq5}
\label{sec:exp:rq6}
\todo{RQ4: expert taxonomies and LLM baselines (zero-shot, chunked RAG, RAPTOR,
TnT-LLM) on the same corpus; coverage, new dimensions accepted by experts,
verification rate. RQ5: ablation rows for grounding strategy, model and prompt
variants beside the manifest result of \cref{sec:exp:rq1}. RQ6: macro-F$_1$ of
classifying Web-sourced token descriptions with and without the ontology, with
one failure case.}

\subsection{Cost}
\label{sec:exp:cost}
Task~A is one call of 302 seconds with 85,729 output tokens, 2,542 uncached
input tokens and 235,298 tokens read from cache: at most \$0.33 at list price
with cached tokens charged in full. Task~B is \BCalls\ calls in \BSeconds\
seconds with \BOutTokens\ output tokens and \BCacheTokens\ million cached
tokens read, at most \$11.7 before the cache-read discount. \todo{Exact billing figures; baseline and frontier-model cost.}


\paragraph{Ethical considerations.}
All sources are public legal instruments, guidance and papers, and the tagged
content is the scientific and patent part of DLT-Corpus, not its social-media
part; no personal data or human participants are involved. An induced legal
taxonomy can be misread as legal classification, so every node carries its
citations for a reviewer to check before relying on it.

\section{Conclusion}
\label{sec:conclusion}

We proposed \sys, which induces the classification dimensions of crypto-asset
law and research from definitional text with a citation for every element,
orders them into a DAG by one cited is-a query per node, grows the DAG from
tagged Web content by measured information gain and indexes that content on
it. It produced 93 cited term types with 317 values and \BEdges\ is-a
edges; \NIterWord\ iterations tagged the whole
on-topic DLT-Corpus of \AllItems\ items, verified \EvidMin--\EvidMax\% of tag
evidence and added \AllAddedTypes\ dimensions for \$\AllCost. Against
zero-shot, chunked RAG, RAPTOR and TnT-LLM pipelines over the same corpus and
model, the whole-corpus catalogue scores highest downstream, and without the
corpus the same model recovers a quarter of the expert key; removing the
citation requirement alone cuts the catalogue from 93 to 39 term types. On
an independent task, site categories for Web token descriptions, the
ontology helps as learned features, 15 points of macro-F$_1$ over TF-IDF,
and not as prompt context. Expert validation and regulatory batches are left to the next version.


%% Acknowledgments are omitted in anonymous mode.
\ifcameraready
\begin{acks}
TODO: funding, data providers, reviewers.
\end{acks}
\fi

%% ============================================================================
%%  REFERENCES  (do not count towards the 8-page body; total cap 12 pages)
%% ============================================================================
\bibliographystyle{ACM-Reference-Format}
\bibliography{references}

%% ============================================================================
%%  APPENDIX  (optional; reviewers are not required to read it)
%% ============================================================================
\appendix
%% ============================================================================
%%  APPENDIX -- optional. Total paper (body + refs + appendix) must be <= 12 pp.
%%  Reviewers are NOT required to read this; nothing essential belongs here.
%% ============================================================================

\section{Implementation Details}
\label{app:impl}
\todo{Full prompt templates (system instructions, per-term query, JSON output
schema); model identifiers and dates; decoding parameters; cache configuration;
total token usage and cost.}

\section{Corpus Construction Details}
\label{app:corpus}
\todo{Per-instrument section ranges / regexes used for extraction; cleaning
rules; known PDF extraction artefacts.}

\section{Annotation Guidelines}
\label{app:annotation}
\todo{Instructions given to expert annotators; examples; adjudication protocol.}

\section{Additional Results}
\label{app:results}
\todo{Per-jurisdiction and per-term breakdowns; full ablation tables;
significance tests.}

\section{Taxonomy Excerpt}
\label{app:taxonomy}
\begin{figure}[h]
  \centering
  \todo{Rendered subtree of the taxonomy for \term{stablecoin} with
  jurisdiction-indexed leaves.}
  \caption{Excerpt of the induced taxonomy. \todo{caption}}
  \label{fig:taxonomy}
\end{figure}


\end{document}
